% ECONOMICS_PROBLEM: {"id": "C73-391", "title": "Threshold-constrained infinite-horizon equilibrium realizability", "jel_code": "C73", "source_id": "OP-0160"}
% !TeX program = xelatex
\documentclass[11pt,a4paper]{article}
\usepackage[margin=23mm]{geometry}
\usepackage{fontspec}
\setmainfont{Latin Modern Roman}
\usepackage{amsmath,amssymb,mathtools}
\usepackage{xcolor,enumitem,booktabs,longtable,array,xurl,hyperref}
\definecolor{muted}{HTML}{53636C}
\hypersetup{colorlinks=true,urlcolor=blue,linkcolor=blue}
\setlength{\parindent}{0pt}
\setlength{\parskip}{5pt}
\newcommand{\R}{\mathbb R}
\newcommand{\E}{\mathbb E}
\newcommand{\N}{\mathbb N}
\newcommand{\ind}{\mathbf 1}
\newcommand{\Var}{\operatorname{Var}}
\newcommand{\Cov}{\operatorname{Cov}}
\newcommand{\supp}{\operatorname{supp}}
\DeclareMathOperator*{\argmin}{arg\,min}
\DeclareMathOperator*{\argmax}{arg\,max}
\newcommand{\entrymeta}[3]{{\sffamily\small\color{muted}\textbf{\texttt{#1}}\quad|\quad #2\par #3\par}}
\newcommand{\Problemref}[1]{\texttt{#1}}
\newcommand{\JELref}[2]{\texttt{#2} (JEL \texttt{#1})}
\begin{document}
\section*{Threshold-constrained infinite-horizon equilibrium realizability}
\textbf{Problem ID:} \texttt{C73-391}\quad \textbf{JEL:} \texttt{C73}\par
% BEGIN ECONOMICS STATEMENT
\label{problem:OP-0160}
% GAME_THEORY_PROBLEM: {"local_id": "GT-030", "original_number": 30, "kind": "H", "entry_kind": "statement_only", "published": false, "review_required": true, "category": "game_theory"}
\label{game:GT-030}
{\sffamily\small\color{muted}
{\sffamily\small\color{muted}\textbf{GT-030} | Stochastic games and computational complexity\par H/U | General NE decision is undecidable; SPE and restricted variants require separate status checks\par}
\par}
\subsection*{Definitions and mathematical statement}
An input consists of a finite concurrent stochastic game with an explicitly listed rational transition table, initial state $s_0$, target sets $F_i$, and rational thresholds $t_i\in[0,1]$, all encoded in binary. Under complete observation of states and joint actions, allow arbitrary behavioral strategies. Write $u_i(\sigma)=\Pr_{s_0,\sigma}(\exists q:\ s_q\in F_i)$. The constrained Nash language is
\[
 \mathcal L_{\mathrm{NE}}=
 \{(G,s_0,(F_i),(t_i)):\exists\sigma\text{ Nash},\ 
 u_i(\sigma)\geq t_i\text{ for every }i\}.
\]
This general decision problem is already undecidable. [S2], Theorem 4.9, gives non-recursive-enumerability even for ten-player simple stochastic multiplayer games: it is not recursively enumerable to decide whether a Nash equilibrium exists in which a designated player wins almost surely. These turn-based games with absorbing reachability outcomes are contained in the displayed class; the special threshold vector is $t_0=1$, $t_i=0$ for all other players.

The corresponding language $\mathcal L_{\mathrm{SPE}}$ replaces ``Nash'' by equilibrium after every feasible full history, with prefix target achievements retained. Its status is not inferred from the NE theorem.
\subsection*{Short English statement}
The unrestricted payoff-threshold NE problem is undecidable; meaningful remaining classification questions must restrict the game or strategy class.
\subsection*{Variants and scope boundaries}
Fixing two players, allowing only finite-memory or stationary equilibria, requiring approximation, or asking unconstrained existence gives different problems. Verification of an input strategy profile also differs from realizability. A decision question about existence of arbitrary infinite-history strategies does not assume that equilibria admit finite encodings.
\subsection*{Source relationship and status}
The original catalog overextends [S1]'s discussion: that source concerns unconstrained existence, including the two-player question, rather than establishing that unrestricted threshold-constrained NE decision remains open. [S2] corrects this attribution. No unresolved complexity claim is made here for the general NE language, and no precise source-verified classification is claimed for the separately defined SPE language.
\subsection*{References}
\begingroup\footnotesize\raggedright
\textbf{[S1]} Senthil Rajasekaran and Moshe Y. Vardi, \emph{Verifying Equilibria in Finite-Horizon Probabilistic Concurrent Game Systems}, arXiv:2503.14690v7 (2026). Locator: Section 1, discussion distinguishing finite-horizon existence, infinite-horizon existence, and verification. \url{https://arxiv.org/html/2503.14690v7}

\textbf{[S2]} Michael Ummels and Dominik Wojtczak, \emph{The Complexity of Nash Equilibria in Stochastic Multiplayer Games}, Logical Methods in Computer Science 7(3:20) (2011), pp. 1--45. Locator: Theorem 4.9, p. 27, and its reduction to equilibria with player 0 winning almost surely. \url{https://arxiv.org/pdf/1109.4017}
\par\endgroup

\subsection*{Common conventions: Source mathematical conventions}

\textbf{Finite games.} For a finite set $A$, $\Delta(A)$ is its probability
simplex. Players choose independent mixed strategies in Nash equilibrium;
correlation is allowed only when explicitly part of the solution concept.
In a game with payoff functions $u_i$ and strategy profile $x$, an additive
$\varepsilon$ Nash equilibrium satisfies
\[
 u_i(x)\geq u_i(x_i',x_{-i})-\varepsilon
 \quad\text{for every player $i$ and every admissible deviation $x_i'$}.
\]
For numerical approximation, payoffs are normalized as locally specified.
An $\varepsilon$ payoff residual is distinct from distance at most
$\varepsilon$ to the set of exact equilibria. Point convergence, convergence
of empirical play, and convergence in distance to an equilibrium set are
also distinct.

\textbf{Computational representations.} Unless stated otherwise, a complexity
question takes rational data with binary encoding; polynomial time is measured
in total bit length. A fixed-accuracy polynomial algorithm is not necessarily
polynomial in $1/\varepsilon$, and neither is necessarily polynomial in
$\log(1/\varepsilon)$. Succinct games and value, demand, or sampling oracles
must declare their interfaces. Exact real solutions require an explicit
representation; an unspecified list of arbitrary real numbers is not a
finite computational output. A claimed algorithmic barrier is meaningful
only for its specified model and reductions.

\textbf{Dynamic games.} Histories, information, observation, and allowable
randomization are part of the model. A stationary strategy depends only on
the current state; a positional strategy in a deterministic graph game
chooses one action at each controlled vertex. Nash equilibrium at an initial
state and equilibrium after every history are different requirements.
An expectation of a pathwise $\liminf$ payoff, a $\liminf$ of expected averages,
a limiting expected average, and a uniform finite-horizon equilibrium are
different criteria. None is silently substituted for another.

\textbf{Allocation and mechanisms.} A complete allocation assigns every item
exactly once unless otherwise stated. Zero-valued goods, negative values,
capacity constraints, feasibility, lotteries, and copies of goods affect
fairness definitions and are specified locally. Dominant-strategy incentive
compatibility, Bayesian incentive compatibility, universal truthfulness,
and truthfulness in expectation impose different inequalities. Whenever
an objective ratio has zero denominator, its convention is stated or those
instances are excluded.

\textbf{Infinite-dimensional models.} Expectations, integrals, stochastic
equations, and optimization objectives require their local regularity,
integrability, filtrations, and admissible domains. A backward--forward
mean-field system is a boundary-value problem; it is not an initial-value
semiflow without an additional construction. Long-time, large-population,
and vanishing-noise limits require an order of limits and a convergence
topology. If a minimum need not be attained, the objective is its infimum
or supremum and the approximation requirement is stated separately.
% END ECONOMICS STATEMENT
\end{document}
